
\documentclass{article}
\usepackage{colortbl}
\usepackage{makecell}
\usepackage{multirow}
\usepackage{supertabular}

\begin{document}

\newcounter{utterance}

\centering \large Interaction Transcript for game `eqbench', experiment `eq\_bench', episode 11 with reflexion{-}llama{-}3.1{-}8b{-}instruct.
\vspace{24pt}

{ \footnotesize  \setcounter{utterance}{1}
\setlength{\tabcolsep}{0pt}
\begin{supertabular}{c@{$\;$}|p{.15\linewidth}@{}p{.15\linewidth}p{.15\linewidth}p{.15\linewidth}p{.15\linewidth}p{.15\linewidth}}
    \# & \multicolumn{2}{c}{Player} && \multicolumn{2}{c}{Game Master} \\
    \hline

    \theutterance \stepcounter{utterance}  
    & & & \multicolumn{4}{p{0.6\linewidth}}{
        \cellcolor[rgb]{0.9,0.9,0.9}{
            \makecell[{{p{\linewidth}}}]{
                \texttt{\tiny{[P1$\langle$GM]}}
                \texttt{You have attempted this same task before and did not succeed. These are the plans you wrote after those attempts.} \\
\texttt{Trial \#1: New plan: I failed to accurately predict Harold's emotional response because I didn't consider the context of the conversation. Harold's initial statement was dismissive and somewhat condescending, which likely triggered Lorelei's passionate response. However, Harold's tone and language didn't convey a sense of admiration or disbelief. Instead, he seemed to be questioning Lorelei's perspective, which would likely lead to irritation. Next time, I will pay closer attention to the emotional tone and language used by the characters in the dialogue to make a more accurate prediction. I will also consider the context of the conversation and the relationships between the characters to better understand their emotional responses.} \\
\\ 
\texttt{Trial \#2: New plan: I failed to accurately predict Harold's emotional response because I didn't consider the context of the conversation and the relationships between the characters. In Trial \#1, I realized that Harold's tone and language were dismissive and condescending, which likely triggered Lorelei's passionate response. However, I didn't fully account for the fact that Harold's initial statement was a challenge to Lorelei's perspective, which would likely lead to irritation.} \\
\\ 
\texttt{To improve my prediction, I will pay closer attention to the emotional tone and language used by the characters in the dialogue, as well as the context of the conversation and the relationships between the characters. I will also consider the following specific actions:} \\
\\ 
\texttt{1. Identify the emotional tone of the initial statement: In this case, Harold's initial statement was dismissive and condescending, which set the tone for the rest of the conversation.} \\
\texttt{2. Analyze the character's relationships: Harold and Lorelei seem to have a close relationship, which may influence their emotional responses to each other.} \\
\texttt{3. Consider the context of the conversation: The conversation is about a contentious election, which may be a sensitive topic for both characters.} \\
\texttt{4. Look for emotional cues: I will pay closer attention to emotional cues such as body language, tone of voice, and language used by the characters to better understand their emotional responses.} \\
\\ 
\texttt{By taking these actions, I hope to improve my prediction of Harold's emotional response and achieve a better outcome in the future.} \\
\\ 
\texttt{Trial \#3: New plan: I failed to accurately predict Harold's emotional response because I didn't consider the context of the conversation and the relationships between the characters. In my previous attempts, I realized that Harold's tone and language were dismissive and condescending, which likely triggered Lorelei's passionate response. However, I didn't fully account for the fact that Harold's initial statement was a challenge to Lorelei's perspective, which would likely lead to irritation.} \\
\\ 
\texttt{To improve my prediction, I will pay closer attention to the emotional tone and language used by the characters in the dialogue, as well as the context of the conversation and the relationships between the characters. I will also consider the following specific actions:} \\
\\ 
\texttt{1. Identify the emotional tone of the initial statement: In this case, Harold's initial statement was dismissive and condescending, which set the tone for the rest of the conversation.} \\
\texttt{2. Analyze the character's relationships: Harold and Lorelei seem to have a close relationship, which may influence their emotional responses to each other.} \\
\texttt{3. Consider the context of the conversation: The conversation is about a contentious election, which may be a sensitive topic for both characters.} \\
\texttt{4. Look for emotional cues: I will pay closer attention to emotional cues such as body language, tone of voice, and language used by the characters to better understand their emotional responses.} \\
\\ 
\texttt{To further improve my prediction, I will also consider the following additional actions:} \\
\\ 
\texttt{5. Identify the character's motivations: What are Harold's motivations for making this statement? Is he trying to provoke a reaction from Lorelei or is he genuinely trying to understand her perspective?} \\
\texttt{6. Consider the power dynamics: Who has the power in this conversation? Is Harold trying to assert his dominance or is Lorelei trying to assert her own perspective?} \\
\texttt{7. Look for emotional escalation: How does the conversation escalate emotionally? Does Harold's statement lead to a more intense or heated conversation?} \\
\\ 
\texttt{By taking these actions, I hope to improve my prediction of Harold's emotional response and achieve a better outcome in the future.} \\
\\ 
\\ 
\\ 
\texttt{Your task is to predict the likely emotional responses of a character in this dialogue:} \\
\\ 
\texttt{Harold: I don't understand why you're so worked up about this election. It's not like the world is going to end if your candidate doesn't win.} \\
\texttt{Lorelei: That's not the point, Harold. It's about what this election represents. It's about the future of our planet, our society.} \\
\texttt{Harold: And you think your candidate has all the answers? } \\
\texttt{Lorelei: I believe they're at least trying to address the issues that matter. Climate change, social justice, healthcare...} \\
\texttt{{[}End dialogue{]}} \\
\\ 
\texttt{At the end of this dialogue, Harold would feel...} \\
\texttt{Disbelief} \\
\texttt{Irritation} \\
\texttt{Admiration} \\
\texttt{Indifference} \\
\\ 
\texttt{Give each of these possible emotions a score from 0{-}10 for the relative intensity that they are likely to be feeling each.} \\
\\ 
\texttt{You must output in the following format, including headings (of course, you should give your own scores), with no additional commentary:} \\
\\ 
\texttt{Disbelief: <score>} \\
\texttt{Irritation: <score>} \\
\texttt{Admiration: <score>} \\
\texttt{Indifference: <score>} \\
\\ 
\\ 
\texttt{{[}End of answer{]}} \\
\\ 
\texttt{Remember: zero is a valid score, meaning they are likely not feeling that emotion. You must score at least one emotion > 0.} \\
\\ 
\texttt{Your answer:} \\
            }
        }
    }
     \\ \\

    \theutterance \stepcounter{utterance}  
    & \multicolumn{4}{p{0.6\linewidth}}{
        \cellcolor[rgb]{0.561,0.737,0.561}{
            \makecell[{{p{\linewidth}}}]{
                \texttt{\tiny{[P1$\rangle$GM]}}
                \texttt{Disbelief: 0} \\
\texttt{Irritation: 8} \\
\texttt{Admiration: 0} \\
\texttt{Indifference: 2} \\
            }
        }
    }
    & & \\ \\

    \theutterance \stepcounter{utterance}  
    & & & \multicolumn{2}{p{0.3\linewidth}}{
        \cellcolor[rgb]{0.9,0.9,0.9}{
            \makecell[{{p{\linewidth}}}]{
                \texttt{\tiny{[GM$|$GM]}}
                \texttt{Disbelief: 0} \\
\texttt{Irritation: 8} \\
\texttt{Admiration: 0} \\
\texttt{Indifference: 2} \\
            }
        }
    }
    & & \\ \\

    \theutterance \stepcounter{utterance}  
    & & & \multicolumn{2}{p{0.3\linewidth}}{
        \cellcolor[rgb]{0.9,0.9,0.9}{
            \makecell[{{p{\linewidth}}}]{
                \texttt{\tiny{[GM$|$GM]}}
                \texttt{Disbelief: 4} \\
\texttt{Irritation: 6} \\
\texttt{Admiration: 0} \\
\texttt{Indifference: 0} \\
            }
        }
    }
    & & \\ \\

    \theutterance \stepcounter{utterance}  
    & & & \multicolumn{2}{p{0.3\linewidth}}{
        \cellcolor[rgb]{0.9,0.9,0.9}{
            \makecell[{{p{\linewidth}}}]{
                \texttt{\tiny{[GM$|$GM]}}
                \texttt{game\_result = LOSE} \\
            }
        }
    }
    & & \\ \\

    \theutterance \stepcounter{utterance}  
    & & \multicolumn{4}{p{0.6\linewidth}}{
        \cellcolor[rgb]{0.95,0.95,0.95}{
            \makecell[{{p{\linewidth}}}]{
                \texttt{\tiny{[P1$\leftrightarrow$P]}}
                \texttt{Post{-}episode Reflexion (not sent to GM):} \\
\texttt{New plan: I failed to accurately predict Harold's emotional response because I didn't consider the context of the conversation and the relationships between the characters. In my previous attempts, I realized that Harold's tone and language were dismissive and condescending, which likely triggered Lorelei's passionate response. However, I didn't fully account for the fact that Harold's initial statement was a challenge to Lorelei's perspective, which would likely lead to irritation.} \\
\\ 
\texttt{To improve my prediction, I will pay closer attention to the emotional tone and language used by the characters in the dialogue, as well as the context of the conversation and the relationships between the characters. I will also consider the following specific actions:} \\
\\ 
\texttt{1. Identify the emotional tone of the initial statement: In this case, Harold's initial statement was dismissive and condescending, which set the tone for the rest of the conversation.} \\
\texttt{2. Analyze the character's relationships: Harold and Lorelei seem to have a close relationship, which may influence their emotional responses to each other.} \\
\texttt{3. Consider the context of the conversation: The conversation is about a contentious election, which may be a sensitive topic for both characters.} \\
\texttt{4. Look for emotional cues: I will pay closer attention to emotional cues such as body language, tone of voice, and language used by the characters to better understand their emotional responses.} \\
\texttt{5. Identify the character's motivations: What are Harold's motivations for making this statement? Is he trying to provoke a reaction from Lorelei or is he genuinely trying to understand her perspective?} \\
\texttt{6. Consider the power dynamics: Who has the power in this conversation? Is Harold trying to assert his dominance or is Lorelei trying to assert her own perspective?} \\
\texttt{7. Look for emotional escalation: How does the conversation escalate emotionally? Does Harold's statement lead to a more intense or heated conversation?} \\
\\ 
\texttt{Additionally, I will also consider the following new actions:} \\
\\ 
\texttt{8. Identify the character's emotional triggers: What specific topics or phrases trigger strong emotions in Harold and Lorelei? Understanding these triggers can help me better predict their emotional responses.} \\
\texttt{9. Consider the character's past experiences: Have Harold and Lorelei had past experiences that may influence their emotional responses to this conversation? Understanding their past experiences can help me better understand their motivations and emotional triggers.} \\
\texttt{10. Look for inconsistencies: Are there any inconsistencies in Harold's or Lorelei's statements or actions that may indicate a deeper emotional response? Identifying these inconsistencies can help me better understand their emotional responses.} \\
\\ 
\texttt{By taking these actions, I hope to improve my prediction of Harold's emotional response and achieve a better outcome in the future.} \\
            }
        }
    }
    & \\ \\

\end{supertabular}
}

\end{document}
